
\documentclass{article}
\usepackage{colortbl}
\usepackage{makecell}
\usepackage{multirow}
\usepackage{supertabular}

\begin{document}

\newcounter{utterance}

\centering \large Interaction Transcript for game `wordle\_withcritic\_noexp', experiment `medium\_frequency\_words\_clue\_with\_critic', episode 18 with reflexion{-}Qwen3.5{-}9B.
\vspace{24pt}

{ \footnotesize  \setcounter{utterance}{1}
\setlength{\tabcolsep}{0pt}
\begin{supertabular}{c@{$\;$}|p{.15\linewidth}@{}p{.15\linewidth}p{.15\linewidth}p{.15\linewidth}p{.15\linewidth}p{.15\linewidth}}
    \# & $\;$A & \multicolumn{4}{c}{Game Master} & $\;\:$B\\
    \hline

    \theutterance \stepcounter{utterance}  
    & & \multicolumn{4}{p{0.6\linewidth}}{
        \cellcolor[rgb]{0.9,0.9,0.9}{
            \makecell[{{p{\linewidth}}}]{
                \texttt{\tiny{[P1$\langle$GM]}}
                \texttt{You have attempted this same task before and did not succeed. These are the plans you wrote after those attempts.} \\
\texttt{Trial \#1: New plan: I failed because I repeatedly guessed "judge" despite receiving negative feedback (j, d, e are red; g is yellow) and a "no" agreement, ignoring the clue "court activity." Next time, I will immediately incorporate the clue to select a word related to sports or legal proceedings (e.g., "match", "trial", "bench") and ensure I do not reuse letters marked as red in previous feedback. I will also pay attention to the "guess\_agreement: no" signal as a prompt to change my guess rather than repeating the same one.} \\
\\ 
\texttt{Trial \#2: New plan: I failed because I ignored the "guess\_agreement: no" signal after my first guess and stuck with "judge" despite the feedback showing 'j', 'd', and 'e' are incorrect (red) and 'g' is misplaced (yellow). The clue "court activity" strongly suggests a legal term, but "judge" was ruled out. Next time, upon receiving "guess\_agreement: no" or negative feedback, I will immediately pivot to a different legal term that excludes red letters (j, d, e) and avoids the misplaced letter (g). I will prioritize words like "verdict," "plaintiff," "defendant," or "arrest" that fit the clue and respect the letter constraints.} \\
\\ 
\texttt{Trial \#3: New plan: I failed because I ignored the explicit instruction to use the clue "court activity" and the negative feedback from my previous guess ("judge" was rejected with 'j', 'd', 'e' red and 'g' yellow). Instead of pivoting to a valid legal term that excludes those letters, I likely repeated a guess or provided an invalid response format. Next time, upon receiving "guess\_agreement: no" and negative feedback, I will immediately analyze the clue "court activity" to select a new legal term (e.g., "arrest," "trial," "bench," or "jury") that strictly excludes all red letters (j, d, e) and the misplaced letter (g), ensuring my response contains only the required "guess:" tag and a valid five{-}letter word.} \\
\\ 
\\ 
\texttt{You are a language wizard who likes to guess words by using the given rules.} \\
\\ 
\texttt{Welcome to Wordle! You have six attempts to guess the target word, a valid English word of five lowercase letters (a{-}z). Please use the tag "guess:" to provide your guess.} \\
\\ 
\texttt{To help you make an informed guess, you will receive a clue for the word, such as} \\
\texttt{clue: "snowy white"} \\
\\ 
\texttt{Here is an example guess based on the clue: } \\
\texttt{guess: apple} \\
\\ 
\texttt{I will then indicate whether I agree or disagree with your guess, but agreeing with a guess does not confirm its correctness. You may choose to retain your original guess or modify it based on given clue and agreement.} \\
\\ 
\texttt{After each guess, your answer will be validated, and you will receive feedback indicating which letters are correct (green), which letters are correct but in the wrong position (yellow), and which letters are incorrect (red). This feedback can be useful in determining which letters to include or exclude in your next guess.} \\
\\ 
\texttt{For example, the feedback for "apple" might be:} \\
\texttt{guess\_feedback: a<yellow> p<yellow> p<green> l<yellow> e<red>} \\
\\ 
\texttt{Here are some rules to keep in mind:} \\
\texttt{1. Please do not generate any guess feedback or clue on your own} \\
\texttt{2. Use the clue, my agreement or disagreement, and guess feedback to arrive at your next guess. } \\
\\ 
\texttt{Do not generate any other text.} \\
\\ 
\texttt{Let's begin with your first guess.} \\
\\ 
\texttt{clue: court activity} \\
            }
        }
    }
    & \\ \\

    \theutterance \stepcounter{utterance}  
    & \multicolumn{4}{p{0.6\linewidth}}{
        \cellcolor[rgb]{0.9,0.9,0.9}{
            \makecell[{{p{\linewidth}}}]{
                \texttt{\tiny{[P1$\rangle$GM]}}
                \texttt{guess: trial} \\
            }
        }
    }
    & & \\ \\

    \theutterance \stepcounter{utterance}  
    & & \multicolumn{4}{p{0.6\linewidth}}{
        \cellcolor[rgb]{0.9,0.9,0.9}{
            \makecell[{{p{\linewidth}}}]{
                \texttt{\tiny{[GM$\rangle$P2]}}
                \texttt{You have attempted this same task before and did not succeed. These are the plans you wrote after those attempts.} \\
\texttt{Trial \#1: New plan: I failed because I repeatedly answered "no" without ever analyzing the clue "court activity" or the guess "judge" to determine if they actually matched. The task requires me to agree or disagree based on the semantic relationship between the clue and the guess, not to generate a random sequence of rejections. Since "judge" is a person involved in court activities, the correct answer is "yes". Next time, I will immediately evaluate the semantic link between the clue and the provided guess and output "agreement: yes" if they match, adhering strictly to the required format.} \\
\\ 
\texttt{Trial \#2: New plan: I failed because I incorrectly agreed to the guess "judge" on the first turn, likely due to a premature assumption that the semantic link was strong enough without verifying the specific constraints of the game loop or the exact definition of "court activity" versus "judge." The game ended in a loss because I stopped the interaction too early or failed to recognize that the task required a specific sequence of validation steps that I missed. Next time, I will strictly analyze the clue "court activity" against the guess "judge" to ensure the semantic match is definitive before outputting "agreement: yes." If the connection is ambiguous or if the game rules imply a need for further verification (like checking for feedback availability), I will output "agreement: no" to allow the game to proceed with a new guess, rather than forcing an agreement that might be premature. I will also ensure I do not repeat the same "yes" response if the game state suggests a need for re{-}evaluation.} \\
\\ 
\\ 
\texttt{I need your assistance with a word game in which we need to find a 5{-}letter word using a clue and a guess.} \\
\\ 
\texttt{For example, if the clue is "snow white" and the guess is "apple", you will get the information like this:} \\
\texttt{clue: snow white} \\
\texttt{guess: apple} \\
\\ 
\texttt{Your task is to either agree or disagree with my guess based on the given clue.} \\
\\ 
\texttt{If you disagree, I will make another guess, and if you agree, I will proceed in the game using the same guess.} \\
\\ 
\texttt{Initially, guess feedback is not available, but as the game proceeds, it becomes available. The guess feedback is not for current guess, but rather an indication of what letters can be present in the current guess. A letter highlighted in green indicates that it is a correct letter in the correct position, while yellow indicates that it is a correct letter in the wrong position, and red indicates a wrong letter. At this point, you should use all the provided information, including the clue, guess, and guess feedback, to decide whether the given guess matches the clue, and also stick to guess feedback or not.} \\
\\ 
\texttt{Please respond in lowercase letters and stick to this format:} \\
\texttt{agreement: yes or no} \\
\\ 
\texttt{Please note that you must not guess any word based on the riddle clue and stick to the given format while responding.} \\
\\ 
\texttt{Do not generate any other text.} \\
\\ 
\texttt{Let's begin!} \\
\\ 
\texttt{clue: court activity} \\
\texttt{guess: trial} \\
            }
        }
    }
    & \\ \\

    \theutterance \stepcounter{utterance}  
    & & & \multicolumn{4}{p{0.6\linewidth}}{
        \cellcolor[rgb]{0.9,0.9,0.9}{
            \makecell[{{p{\linewidth}}}]{
                \texttt{\tiny{[GM$\langle$P2]}}
                \texttt{agreement: no} \\
            }
        }
    }
     \\ \\

    \theutterance \stepcounter{utterance}  
    & & \multicolumn{4}{p{0.6\linewidth}}{
        \cellcolor[rgb]{0.9,0.9,0.9}{
            \makecell[{{p{\linewidth}}}]{
                \texttt{\tiny{[P1$\langle$GM]}}
                \texttt{clue: court activity} \\
\texttt{guess\_agreement: no} \\
\texttt{Provide your response only in this format.} \\
\texttt{guess: word} \\
            }
        }
    }
    & \\ \\

    \theutterance \stepcounter{utterance}  
    & \multicolumn{4}{p{0.6\linewidth}}{
        \cellcolor[rgb]{0.9,0.9,0.9}{
            \makecell[{{p{\linewidth}}}]{
                \texttt{\tiny{[P1$\rangle$GM]}}
                \texttt{guess: verdict} \\
            }
        }
    }
    & & \\ \\

    \theutterance \stepcounter{utterance}  
    & & & \multicolumn{2}{p{0.3\linewidth}}{
        \cellcolor[rgb]{0.9,0.9,0.9}{
            \makecell[{{p{\linewidth}}}]{
                \texttt{\tiny{[GM$|$GM]}}
                \texttt{Error: The length of the guessed word is not 5.} \\
            }
        }
    }
    & & \\ \\

    \theutterance \stepcounter{utterance}  
    & & & \multicolumn{2}{p{0.3\linewidth}}{
        \cellcolor[rgb]{0.9,0.9,0.9}{
            \makecell[{{p{\linewidth}}}]{
                \texttt{\tiny{[GM$|$GM]}}
                \texttt{game\_result = ABORT} \\
            }
        }
    }
    & & \\ \\

    \theutterance \stepcounter{utterance}  
    & \multicolumn{4}{p{0.6\linewidth}}{
        \cellcolor[rgb]{0.95,0.95,0.95}{
            \makecell[{{p{\linewidth}}}]{
                \texttt{\tiny{[P2$\rightarrow$P1]}}
                \texttt{Post{-}episode Reflexion (not sent to GM):} \\
\texttt{New plan: I failed because I ignored the explicit instruction to use the clue "court activity" and the negative feedback from my previous guess ("judge" was rejected with 'j', 'd', 'e' red and 'g' yellow). Instead of pivoting to a valid legal term that excludes those letters, I likely repeated a guess or provided an invalid response format. Next time, upon receiving "guess\_agreement: no" and negative feedback, I will immediately analyze the clue "court activity" to select a new legal term (e.g., "arrest," "trial," "bench," or "jury") that strictly excludes all red letters (j, d, e) and the misplaced letter (g), ensuring my response contains only the required "guess:" tag and a valid five{-}letter word.} \\
            }
        }
    }
    & & \\ \\

    \theutterance \stepcounter{utterance}  
    & & & \multicolumn{4}{p{0.6\linewidth}}{
        \cellcolor[rgb]{0.95,0.95,0.95}{
            \makecell[{{p{\linewidth}}}]{
                \texttt{\tiny{[P1$\rightarrow$P2]}}
                \texttt{Post{-}episode Reflexion (not sent to GM):} \\
\texttt{New plan: I failed because I output "agreement: no" without verifying if the guess "trial" actually matches the clue "court activity." The word "trial" is a direct synonym for a court activity, so the correct semantic evaluation is that they match. I should have analyzed the relationship between the clue and the guess first, determined that "trial" fits "court activity" perfectly, and then output "agreement: yes" in the exact required lowercase format. Next time, I will immediately check if the guess is a valid instance of the clue; if it is, I will output "agreement: yes", and only if it is not will I output "agreement: no".} \\
            }
        }
    }
     \\ \\

\end{supertabular}
}

\end{document}
